\documentclass[a4paper,11pt]{article}
\usepackage[T1]{fontenc}
\usepackage[utf8]{inputenc}

\setlength{\evensidemargin}{0pt} % Marge gauche sur pages paires
\setlength{\marginparwidth}{54pt} % Largeur de note dans la marge
\setlength{\textwidth}{481pt} % Largeur de la zone de texte (17cm)
\setlength{\voffset}{-18pt} % Bon pour DOS
\setlength{\marginparsep}{7pt} % Séparation de la marge
\setlength{\topmargin}{0pt} % Pas de marge en haut
\setlength{\headheight}{13pt} % Haut de page
\setlength{\headsep}{4pt} % Entre le haut de page et le texte
\setlength{\footskip}{27pt} % Bas de page + séparation
\setlength{\textheight}{720pt} % Hauteur de la zone de texte (25cm)
\setlength{\parindent}{0pt}

\usepackage{natbib}
\usepackage{graphicx}
\usepackage{amssymb}
\usepackage{amsmath} % Required for the align* environment
\usepackage{float}
\usepackage{pgfplots}
\usepackage{algorithm}
\usepackage{algpseudocode}
\pgfplotsset{compat=1.18}


\begin{document}
\section{Föreläsning 1}
Nummerisk metod -- Algoritm som löser problem approximativt:
\begin{itemize} 
    \item Ekvationslösning
    \begin{itemize}
        \item För funktioner av en variabel
            \[
            f(x) = 0
            \]
        \item För system av ekvationer
            \[
            \begin{cases}
            f_1(x_1, x_2, \ldots, x_n) = 0 \\
            f_2(x_1, x_2, \ldots, x_n) = 0 \\
            \vdots \\
            f_n(x_1, x_2, \ldots, x_n) = 0
            \end{cases}
            \]
    \end{itemize}
    \item Kurvanpassning
    \begin{itemize}
        \item Interpolation, minsta kvadrat, etc.
    \end{itemize}
    \item Integration 
    \[
    \int_a^b f(x) dx \approx \sum_{i=1}^n f(x_i) \Delta x
    \]
    \item Derivering
    \[
    f'(x) \approx \frac{f(x + \Delta x) - f(x)}{\Delta x}
    \]
    \item Ordinära differentialekvationer
    \[
    \frac{dy(t)}{dt} = f(t, y(t)), \hspace{10pt} y(0) = y_0
    \]
\end{itemize}

\section{Föreläsning 2}
\subsection{Iterativ metod}
\[
x_0, x_1, x_2, \ldots, x_n \rightarrow x^*, \hspace{10pt} f(x^*) = 0
\]
Exempelvis genom intervall halvering och fixpunkt iteration. Den generlla metoden är:
\begin{enumerate}
    \item Hitta en startgissning $x_0$.
    \item Iterera $x_{n} \rightarrow x_{n+1}$
    \item Upprepa tills $|x_{n+1} - x_n| < \epsilon$.
    \item De vanliga frågorna är om den konvergerar och dess konvergens ordning
\end{enumerate}

\subsection{Intervall halvering}
Antag att $f$ är en kontinuerlig funktion samt:
\begin{itemize}
    \item $a_0 < \sqrt{x^*} < b_0$
    \item $x_0 = \frac{a_0 + b_0}{2}$
    \item Om $f(a_0)$ och $f(x_0)$ har olika tecken 
    \begin{itemize}
        \item Byt till $[a_0, x_0]$
    \end{itemize}
    \item Om $f(a_0)$ och $f(x_0)$ har samma tecken
    \begin{itemize}
        \item Byt till $[x_0, b_0]$
    \end{itemize}
    \item Basically binär sökning för att hitta när $f(x) = 0$
\end{itemize}

\subsubsection{Konvergens ordning för intervall halvering}
\begin{itemize}
    \item Intervall längden halveras varje iteration
    \item Detta gör konvergens ordningen 1, eftersom $|x_{n+1} - x^*| \approx \frac{1}{2} |x_n - x^*|$.
\end{itemize}

\subsection{Fixpunktsiteration}
Givet en funktion $f(x) = 0$ kan den alltid skrivas i fixpunktsform:
\[
x = \Phi(x)
\]
Då gäller:
$$x^* = \Phi(x^*) \iff f(x^*) = 0$$
Metoden är:
\begin{enumerate}
    \item Börja med en startgissning $x_0$.
    \item Iterera $x_{n+1} = \Phi(x_n)$.
    \item Upprepa tills $|x_{n+1} - x_n| < \epsilon$
\end{enumerate}

\subsubsection{Analys av fixpunktsiteration}
$$|x_{n+1} - x^*| = |\Phi'(x_n)||x_n - x^*|$$
Felet avtar med faktorn $|\Phi'(x_n)|$ för varje iteration som kan bevisas:
$$x_{n+1} - x^* = \Phi(x_n) - \Phi(x^*)$$
Medelvärdessatsen säger att:
$$f(b) - f(a) = f'(c)(b-a), \hspace{10pt} a < c < b$$
$$\implies \Phi(x_n) - \Phi(x^*) = \Phi'(c)(x_n - x^*), \hspace{10pt} x_n < c < x^*$$
$$\implies |x_{n+1} - x^*| = |\Phi'(c)||x_n - x^*|$$

Där för gäller det att konvergensen beror på $|\Phi'(c)|$:
\begin{itemize}
    \item Om $|\Phi'(c)| < 1$ så konvergerar metoden
    \item Om $|\Phi'(c)| > 1$ så divergerar metoden
    \item Om $|\Phi'(c)| = 1$ så kan det vara svårt att avgöra konvergensen
    \item Om $|\Phi'(c)| = 0$ fås kvadratisk konvergens
\end{itemize}

Kvadratisk konvergens innebär:
$$|x_{n+1} - x^*| \approx S|x_n - x^*|^2, \hspace{10pt} C > 0$$
För fixpunktsiteration vid kvadratisk konvergens gäller att:
$$|x_{n+1} - x^*| \approx \frac{|\Phi''(c)|}{2}|x_n - x^*|^2$$
Generellt om $|\Phi^{(n)}(c)| = 0$ gäller att:
$$|x_{n+1} - x^*| \approx \frac{|\Phi^{(n)}(c)|}{n!}|x_n - x^*|^n$$  

\newpage
\section{Föreläsning 3}
\subsection{Newtons metod}
Givet $f(x) = 0$ så kan man linjärisera kring punkten $x_n$:
$$f(x) \approx L(x) = f(x_n) + f'(x_n)(x - x_n)$$
Låt $L(x_{n+1}) = 0$
$$\implies 0 = L(x_{n+1}) = f(x_n) + f'(x_n)(x_{n+1} - x_n)$$
$$\implies 0 = f(x_n) + f'(x_n)(x_{n+1} - x_n)$$
$$\implies x_{n+1} = x_n - \frac{f(x_n)}{f'(x_n)}$$

\subsubsection{Analys av Newtons metod}
\begin{itemize}
    \item $\exists f'''(x) \implies \text{Lokal konvergens}$
    \item Lokal konvergens innebär att metoden endast konvergerar om stargissningen är nära nog
    \item $f'(x) = 0$ händer inte i kursen
    \item Bevis nedan:
\end{itemize}
\noindent Givet att fixpunktsiterationer är lokalt konvergenta om $|\Phi(x^*)| < 1$:

$$\Phi(x) = x - \frac{f(x)}{f'(x)}$$
$$\implies \frac{d}{dx}\Phi(x) = 1 - \frac{f'(x)^2 - f(x)f''(x)}{f'(x)^2}$$
$$\implies \Phi'(x^*) = 1 - \frac{f'(x^*)^2 - f(x^*)f''(x^*)}{f'(x^*)^2}$$
$$\implies \Phi'(x^*) = 1 - \frac{f'(x^*)^2}{f'(x^*)^2} + \frac{f(x^*)f''(x^*)}{f'(x^*)^2}$$
$$\implies \Phi'(x^*) = \frac{f(x^*)f''(x^*)}{f'(x^*)^2}$$
$$f(x^*) = 0 \implies \Phi'(x^*) = 0$$
Gäller endast om $f'''(x)$ existerar, vilket ger kvadratisk konvergens

\subsubsection{Newtons metod konvergens}
\begin{itemize}
    \item Newtons metod har kvadratisk konvergens eftersom $|\Phi'(x^*)| = 0$.
    \item Vi får $|x_{n+1} - x^*| \approx \frac{|\Phi''(c)|}{2}|x_n - x^*|^2$
\end{itemize}
\newpage
\subsection{Sekant metoden}
\begin{itemize}
    \item Kräver två startgissningar $x_n$ och $x_{n-1}$.
    \item Basically Newton med approximerad derivata:
    \item Lokalt konvergent om $f'(x^*) \neq 0$
\end{itemize}
Anta:
\[
f'(x_n) \approx \frac{f(x_n) - f(x_{n-1})}{x_n - x_{n-1}}
\]
Då blir iterationsformeln:
\[
x_{n+1} = x_n - f(x_n)\frac{x_n - x_{n-1}}{f(x_n) - f(x_{n-1})}
\]

\subsection{Approximativ newton}
Om $f'(x)$ är odefinierad eller svår att beräkna kan kan approximera $f'(x_n)$ som:
$$f'(x_n) \approx \frac{f(x_n + h) - f(x_n)}{h}, \hspace{10pt} h = 10^{-5} \approx 0$$
Detta ger då:
$$x_{n+1} = x_n - \frac{h \cdot f(x_n)}{f(x_n + h) - f(x_n)}$$

\subsection{Konvergens}
\subsubsection{Linjärkonvergens}
Felet i $x_{n+1}$ är ungefär proportionellt $x_n$
$$|x_{n+1} - x^*| \approx C|x_n - x^*|, \hspace{10pt} C > 0$$
$$\frac{|x_{n+1} - x^*|}{|x_n - x^*|} \approx C$$

\subsubsection{Kvadratisk konvergens}
$$|x_{n+1} - x^*| \approx S|x_n - x^*|^2, \hspace{10pt} C > 0$$
$$\frac{|x_{n+1} - x^*|}{|x_n - x^*|^2} \approx C$$
Antalet korrekta dubblas efter varje iteration

\section{Föreläsning 4}











\end{document}